\documentclass[11pt,a4paper]{article}
\usepackage{fontspec,polyglossia}
\setdefaultlanguage{russian}
\setmainfont{Arial}
\newfontfamily\cyrillicfont{Arial}
\usepackage[margin=19mm,headheight=15pt]{geometry}
\usepackage{amsmath,xcolor,needspace,setspace,fancyhdr,xstring}
\usepackage[hidelinks]{hyperref}
% Числа для презентации и речи. Источники: Paper/main.tex.
% tab:impacts и пояснение после таблицы.
\newcommand{\CompetitorEffect}{-0{,}027}
\newcommand{\CompetitorPercent}{2{,}66\%}
% tab:sdm, tab:robustness.
\newcommand{\SpatialCoefficient}{0{,}474}
\newcommand{\RegionalSpatialCoefficient}{0{,}204}
% Часть II, Выводы. MAPE восстановлена по округленным баллам.
\newcommand{\InitialScore}{90{,}78}
\newcommand{\FinalScore}{90{,}83}
\newcommand{\InitialMape}{4{,}7215\%}
\newcommand{\FinalMape}{4{,}6952\%}
\newcommand{\MapeGain}{0{,}0262}

\definecolor{main}{HTML}{35444D}
\definecolor{accent}{HTML}{4F8DB7}
\setstretch{1.08}
\setlength{\emergencystretch}{1em}
\setlength{\parindent}{0pt}
\setlength{\parskip}{4pt}
\pagestyle{fancy}\fancyhf{}
\lhead{\small Пространственная эконометрика}
\rhead{\small Речь к презентации}
\cfoot{\thepage}
\newcommand{\hint}[1]{\StrSubstitute{#1}{min}{мин}[\russianhint]\textcolor{main}{\small(\russianhint)}}
\newcommand{\scriptopening}[2]{\Needspace{7\baselineskip}\subsection*{Пространственные эффекты в выручке продуктового ритейла}\hint{#1}\par #2\par}
\newcommand{\scriptframe}[3]{\par\vspace{4pt}\noindent\begin{minipage}{\linewidth}\setlength{\parskip}{4pt}\textcolor{accent}{\rule{\linewidth}{0.5pt}}\par\vspace{2pt}\textbf{#1}\par\hint{#2}\par\vspace{2pt} #3\par\end{minipage}\par}
\begin{document}
{\LARGE\bfseries Речь к презентации}\par\vspace{2mm}
\textbf{Регламент:} до 5 минут на часть I и до 5 минут на часть II. Картам и диаграмме даем короткое устное пояснение. Формулы и таблицы целиком не читаем. Тайминг окончательно проверяется на репетиции.

\textbf{Выступающие:} Денис, слайды 1-5; Алексей, 6-12; Богдан, 13-16; Никита, 17-19; Владимир, 20-24. У каждого один непрерывный блок.

\scriptopening{Дугарев Денис; слайд 1; 0.06 min}{Здравствуйте. Представляем исследование пространственных эффектов в выручке магазинов «Пятерочка».}

% 2. Источники: Paper/main.tex; для части II - Vladimir_part/X5_slides/frames.tex.
\scriptframe{Вопрос исследования и план}{Дугарев Денис; слайд 2; 0.23 min}{%
Наш исследовательский вопрос - как выручка магазина связана с конкуренцией и окружением соседних локальных рынков? Работа поделена на эконометрический анализ и улучшение нашего прогноза по итогам соревнования X5 Retail Group.
}

% 4. Источники: Paper/main.tex; для части II - Vladimir_part/X5_slides/frames.tex.
\scriptframe{Данные: локальный рынок и выручка магазина}{Дугарев Денис; слайд 4; 0.70 min}{%
Основой послужил датасет X5 Group: выручка, население и домохозяйства, пеший и автомобильный трафик, конкуренты, плотность сети, пункты выдачи, медицинские учреждения, кассы, площадь и возраст магазинов. Добавили население, зарплату, безработицу и индексы цен Росстата, координаты населенных пунктов GeoNames.

Предположили искусственное преобразование и шум в исходных данных, поэтому усреднили характеристики по рынкам и уточнили население. Сначала хотели агрегировать по субъектам, но их 67, а ДЗ требует минимум 80 объектов панели. Точные координаты магазинов удалены из выгрузки, поэтому выбрали населенные пункты.

В панели используем месячный логарифм выручки на магазин, в срезе - его среднее за 26 месяцев, с чертой сверху.
}

% 5. Источники: Paper/main.tex; для части II - Vladimir_part/X5_slides/frames.tex.
\scriptframe{Локальные рынки в нашей выборке}{Дугарев Денис; слайд 5; 0.12 min}{%
Вот какие локальные рынки представлены в нашей выборке. Размер точки показывает число магазинов. Больше всего рынков в европейской части России. Далее Алексей.
}

% 6. Источники: Paper/main.tex; для части II - Vladimir_part/X5_slides/frames.tex.
\scriptframe{Гипотезы: конкуренция и связи соседних рынков}{Гончаров Алексей; слайд 6; 0.16 min}{%
Гипотезы касаются сходства выручки соседей, конкуренции, характеристик окружения и связи месячных отклонений. Проверили также преимущество пространственной модели перед обычной регрессией.
}

% 7. Источники: Paper/main.tex; для части II - Vladimir_part/X5_slides/frames.tex.
\scriptframe{Матрицы весов: что сравнили и выбрали}{Гончаров Алексей; слайд 7; 0.28 min}{%
Сравнили шесть матриц: экспоненциальную, ближайших соседей, обратных расстояний и их квадратов, пороговую и близость по населению. Выбрали экспоненциальную по правдоподобию SAR и SDM. Пятьдесят километров задают скорость убывания веса, соседей учитываем до трехсот.
}

% 8. Источники: Paper/main.tex; для части II - Vladimir_part/X5_slides/frames.tex.
\scriptframe{Пространственная связь выручки: глобальные индексы}{Гончаров Алексей; слайд 8; 0.12 min}{%
Предварительно видно: высокая выручка обычно соседствует с высокой, а низкая - с низкой.
}

% 9. Источники: Paper/main.tex; для части II - Vladimir_part/X5_slides/frames.tex.
\scriptframe{Где находятся кластеры высокой и низкой выручки}{Гончаров Алексей; слайд 9; 0.12 min}{%
Здесь показаны кластеры высокой и низкой выручки по локальному индексу Морана.
}

% 10. Источники: Paper/main.tex; для части II - Vladimir_part/X5_slides/frames.tex.
\scriptframe{Локальные статистики: Гири и Гетис-Орд}{Гончаров Алексей; слайд 10; 0.12 min}{%
Здесь - карты других локальных статистик. Подробно останавливаться на них не будем.
}

% 11. Источники: Paper/main.tex; для части II - Vladimir_part/X5_slides/frames.tex.
\scriptframe{Выбор пространственной модели}{Гончаров Алексей; слайд 11; 0.31 min}{%
Перейдем сразу к моделям. Использовали подходы от простой модели к сложной и от общей к частной. В модели выделяются выручка соседей, свои и соседние характеристики и зависимость ошибок. Выбрали SDM по LR, AIC, диагностике остатков и робастному Вальду. BIC предпочел SAR.
}

% 12. Источники: Paper/main.tex; для части II - Vladimir_part/X5_slides/frames.tex.
\scriptframe{SDM: конкуренция и региональная среда}{Гончаров Алексей; слайд 12; 0.32 min}{%
Рост среднего числа конкурентов на единицу связан с выручкой на 2,6-2,7 процента ниже в обеих моделях. Эффект собственной сети незначим. После учета субъектов косвенные эффекты потеряли значимость, а пространственный коэффициент снизился с 0,47 до 0,20. Далее Богдан.
}

% 13. Источники: Paper/main.tex; для части II - Vladimir_part/X5_slides/frames.tex.
\scriptframe{Панель: связь месячных отклонений}{Хмельницкий Богдан; слайд 13; 0.28 min}{%
Учли постоянный уровень каждого рынка и общие эффекты месяцев. Сравнили модели с лагом выручки, с пространственной ошибкой и с обоими компонентами. Отдельная FE-проверка убирает общую месячную динамику каждого субъекта; затем проверяем остатки.
}

% 14. Источники: Paper/main.tex; для части II - Vladimir_part/X5_slides/frames.tex.
\scriptframe{Панель: общие шоки имеют значение}{Хмельницкий Богдан; слайд 14; 0.38 min}{%
При отдельном учете пространственные параметры около 0,67. В модели с обоими компонентами сильнее свидетельства зависимости ошибок. После региональных месячных эффектов глобальный индекс Морана остатков стал небольшим, но CD-тест Песарана остался значимым. Остаточная зависимость сохраняется. Стандартные ошибки ML не учитывают возможную связь ошибок между месяцами.
}

% 15. Источники: Paper/main.tex; для части II - Vladimir_part/X5_slides/frames.tex.
\scriptframe{Динамика: инерция и связь во времени}{Хмельницкий Богдан; слайд 15; 0.21 min}{%
Добавили прошлую выручку для учета инерции. Сопоставили МНК, FE и GMM. В пространственной версии выделили сегодняшнюю и прошлую выручку соседей.
}

% 16. Источники: Paper/main.tex; для части II - Vladimir_part/X5_slides/frames.tex.
\scriptframe{Динамика: инерция и ограничения GMM}{Хмельницкий Богдан; слайд 16; 0.27 min}{%
GMM-модели не прошли J-тест. Скорректированная пространственная SAR дала инерцию около 0,83. Ограничение передачи прошлой выручки на пяти процентах не отвергнуто: данные согласуются с одновременными общими шоками, но механизм не доказан. Далее Никита.
}

% 17. Источники: Paper/main.tex; для части II - Vladimir_part/X5_slides/frames.tex.
\scriptframe{GWR: одинаковы ли связи по территории?}{Русских Никита; слайд 17; 0.15 min}{%
GWR оценили по всем характеристикам среза, разрешив коэффициентам меняться по территории. Размер окна подобрала кросс-валидация.
}

% 18. Источники: Paper/main.tex; для части II - Vladimir_part/X5_slides/frames.tex.
\scriptframe{GWR: конкуренция почти везде отрицательна}{Русских Никита; слайд 18; 0.35 min}{%
Для карты выбрали конкуренцию как главный устойчивый результат.
На карте показаны 124 значимые оценки конкурентов, все отрицательные. Для автотрафика значимы 35 оценок, все положительные. По AICc сравнили обычную регрессию, модель с местными коэффициентами и SDM со связями соседей. GWR лучше МНК, но уступает SDM. Зависимость остатков ограничивает надежность локальной значимости.
}

% 19. Источники: Paper/main.tex; для части II - Vladimir_part/X5_slides/frames.tex.
\scriptframe{Часть I: результаты проверки гипотез}{Русских Никита; слайд 19; 0.80 min}{%
\begin{itemize}
\setlength{\itemsep}{2pt}
\setlength{\parsep}{0pt}
\item \textbf{Гипотеза 1 подтверждена:} пространственная связь есть, глобальный индекс Морана 0,58.
\item \textbf{Гипотеза 2 подтверждена:} связь выручки соседей положительна. Пространственный коэффициент 0,47, с эффектами субъектов 0,20.
\item \textbf{Гипотеза 3 подтверждена частично:} рост среднего числа конкурентов на единицу связан с выручкой примерно на 2,7 процента ниже. Прямой эффект своей сети незначим.
\item \textbf{Гипотеза 4 неустойчива:} косвенная связь с населением соседей отрицательна, с автотрафиком и сетью положительна. С эффектами субъектов значимость теряется.
\item \textbf{Гипотеза 5 подтверждена:} SDM лучше МНК по LR-тесту и AIC; для LR $p<0{,}001$.
\item \textbf{Гипотеза 6 подтверждена при предпосылках ML:} пространственная зависимость месячных отклонений положительна, параметры отдельных SAR и SEM около 0,67.
\end{itemize}
Условные связи означают, что мы учитываем включенные в модель характеристики. Однако это не доказывает, что именно конкуренция или соседство вызывают изменения выручки. Далее Владимир.
}

\par\smallskip
\textit{Слайды 20-24 рассказывает Владимир по своему усмотрению.}

\end{document}
